\begin{theorem}[Dyadic interval net finite-net sibling dependency route]
\label{thm:dyadic-interval-net-finite-net-sibling-dependency-route}
Every $\DyadicIntervalNetUp$ finite interval-net refinement consumed by a downstream compact or Real-facing route first displays the sibling rows
$$
\FiniteLebesgueNumberUp,\quad
\CompactMetricUp\ \text{or}\ \CompactIntervalUp,\quad
\DyadicRatCoreUp,\quad
\RealUp .
$$
The finite Lebesgue-number row supplies the mesh threshold, the compact-metric or compact-interval row supplies the bounded interval source, the dyadic row supplies endpoint and radius arithmetic, and the terminal Real row supplies only the seal reached after the finite cells have been listed. Thus the route links \autoref{ch:concrete-instances-finite-lebesgue-number-namecert}, \autoref{ch:concrete-instances-compactmetric-namecert}, the compact-interval surface of \autoref{ch:concrete-instances-bishop-compact-interval-namecert}, \autoref{ch:concrete-instances-dyadicratcore-namecert}, and \autoref{ch:concrete-instances-real-namecert} without exporting an open-cover compactness theorem, selected subcover, quotient stream equality, countable choice, or ambient $\RealUp$ completeness.
\end{theorem}
\begin{proof}
Project the interval-net carrier through the routed chapter body. Its finite cells are already listed before any Real-facing read is available, so a downstream consumer starts from the displayed mesh threshold and compact source rows.

The dyadic arithmetic row then fixes the endpoint and radius checks for each listed cell. This keeps refinement a finite ledger computation rather than an appeal to a host open-cover compactness principle.

Only after those rows have been displayed may the consumer reach the terminal Real seal. Transport, replay, provenance, and local naming preserve that same finite route, excluding selected subcovers, quotient equality, choice, and ambient completeness.
\end{proof}
